% =====================================================================
%  第 1 节：研究问题的提出
% =====================================================================
\section{研究问题的提出}
\label{sec:problem}

\subsection{从一个具体的图论事实说起}
\label{sec:concrete}

设 $G=(V,E)$ 是一个有限简单无向图，$n=|V|$。记
\begin{itemize}
  \item $\alpha(G)$ 为\emph{独立数}：$G$ 中两两不相邻的顶点集的最大大小；
  \item $\omega(G)$ 为\emph{团数}：$G$ 中最大完全子图的大小；
  \item $\chi(G)$ 为\emph{色数}：使相邻顶点异色所需的最少颜色数。
\end{itemize}

这三个量满足一个初等的不等式：
\begin{equation}
  \chi(G) + \alpha(G) \;\le\; n + 1 .
  \label{eq:chi-alpha}
\end{equation}
证明只需一行：取最大独立集 $S$（$|S|=\alpha$），把 $S$ 整体染成一种颜色，其余 $n-\alpha$ 个顶点各染一种颜色，共需 $1+(n-\alpha)=n-\alpha+1$ 种颜色，故 $\chi\le n-\alpha+1$。

不等式 \eqref{eq:chi-alpha} 有两个值得注意的性质。第一，它是\textbf{紧的}：完全图 $K_n$ 上 $\chi=n,\ \alpha=1$，取等号；空图 $\overline{K_n}$ 上 $\chi=1,\ \alpha=n$，也取等号。第二，它的\textbf{证明只需要一行}，但它的\textbf{发现}并不显然——它把一个关于染色的问题（$\chi$）和一个关于独立集的问题（$\alpha$）通过顶点总数 $n$ 联系了起来。

图论中这类\textbf{不变量之间的不等式}数量极多，且其中相当一部分至今未被证明。例如

\begin{conjecture}[Reed, 1998]
\label{conj:reed}
对任意图 $G$，
\begin{equation}
  \chi(G) \;\le\; \ceil{\frac{\omega(G)+\Delta(G)+1}{2}},
  \label{eq:reed}
\end{equation}
其中 $\Delta(G)$ 为最大度。
\end{conjecture}

式 \eqref{eq:reed} 至今未被证明也未被否证。它的两个端点情形分别是已知定理：当 $\omega\le 2$ 时它退化为 Brooks 定理（$\chi\le\Delta+1$，除完全图与奇圈外），当 $\omega=\Delta+1$ 时它退化为平凡界 $\chi\le\Delta+1$。

\subsection{自动猜想：四十年的尝试与同一个瓶颈}
\label{sec:history-brief}

式 \eqref{eq:chi-alpha} 与式 \eqref{eq:reed} 形态相同：都是若干个图不变量的线性组合之间的不等式。一个自然的想法是：\textbf{既然这类不等式的形式如此规整，能否让机器自动地把它们找出来？}

这个想法已经被尝试了将近四十年。

\begin{itemize}
  \item \textbf{Graffiti}（Fajtlowicz，休斯敦大学，1986 年起）。第一个系统性的自动猜想程序。它从一个图数据库中提取不变量数值，用一套人工设计的启发式规则筛出“看起来成立”的不等式，累计提出逾千条猜想，其中若干条后来被数学家证明或否证。
  \item \textbf{Graffiti.pc}（DeLaViña，2000 年代）。Graffiti 的改进版本，调整了启发式以偏向“更可能可证”的猜想。
  \item \textbf{AutoGraphiX}（Caporossi \& Hansen，GERAD，2000 年代）。思路转向：不再从数据中“看”规律，而是用\textbf{优化}主动搜索使某个不等式最可能被违反的图，即主动构造反例；若搜索失败，则把该不等式作为猜想提出。
  \item \textbf{TxGraffiti}（Davila，2020 年代）。引入机器学习与线性规划，从数据中提取系数，并配套给出“极值图”作为猜想的证据。
\end{itemize}

四十年间，方法从启发式规则演进到优化、再到机器学习，但\textbf{有一个瓶颈始终没有变化}：

\begin{keybox}[核心观察]
这些系统的产出\textbf{远多于}人类能够处理的数量。Graffiti 提出逾千条猜想，其中被认真研究过的只是少数；其余的不是错的，而是\textbf{没人知道该不该花时间去证明}。

换句话说，自动猜想系统的困难不在于“生成猜想”——生成从来不是瓶颈——而在于\textbf{从海量候选中判断哪些值得投入}。
\end{keybox}

这个瓶颈有一个明确的技术根源。候选猜想是在\textbf{有限}的图样本上被检验的，而猜想本身断言的是\textbf{对所有图}成立。二者之间的鸿沟意味着：任何候选都可能是“在小图上碰巧成立”的假阳性。当候选数量达到 $10^3$–$10^4$ 量级时，即使每条候选的假阳性率很低，\textbf{期望的假阳性条数}仍然可能远超真实猜想条数。

于是问题变成了一个统计问题：\textbf{如何在大量并行检验中，控制“被接受的猜想里有多少是错的”？}这正是多重检验理论（multiple testing）的核心议题，而据我们所知，上述系统无一系统性地处理过它。

\subsection{本报告研究的问题}
\label{sec:question}

基于以上观察，本报告研究的\textbf{不是}某一个具体的图论猜想（如 Reed 猜想），而是一个方法层面的问题：

\begin{problem}[本课题的研究问题]
\label{prob:main}
能否构建一个智能体，使其在极值图论的不变量不等式中自主地完成“提出猜想—搜寻反例—尝试证明”的循环，并以统计方法控制其产出的假阳性率，使最终保留下来的猜想具备可辩护的可信度？
\end{problem}

这个问题包含三个相互耦合的子问题：

\begin{enumerate}[label=(\roman*), leftmargin=2.4em]
  \item \textbf{猜想生成}：给定一批图及其不变量数值，如何在\textbf{开放的形式空间}中提出候选不等式？（“开放”指：用哪些不变量、取线性还是非线性形式、是否含取整，这些都不预设。）
  \item \textbf{反例搜索}：给定一条候选，如何在\textbf{有限算力}下尽可能快地找到违反它的图？若找不到，如何说明“找不到”是有意义的而非搜索不充分？
  \item \textbf{统计筛选}：给定 $M$ 条候选及各自的检验结果，如何决定保留哪些？如何在候选之间\textbf{分配搜索预算}——把算力花在最有希望或最可疑的候选上？
\end{enumerate}

这三个子问题并非独立：生成策略决定了候选的分布，候选分布决定了筛选的难度，而筛选的反馈又应当指导下一轮生成。这种耦合正是需要一个\textbf{智能体}（而非一段脚本）来承载的原因——见第 \ref{sec:roadmap} 节。

\subsection{为什么选择这个问题}
\label{sec:why}

选择这一问题的理由有三条。

\paragraph{第一，它与统计学的训练直接相关。}
子问题 (iii) 是多重检验与最优实验设计的标准场景：$M$ 条候选对应 $M$ 次检验，需要在控制假发现率的前提下最大化真实发现数；预算分配则是 bandit 问题。这些工具在纯数学背景的研究者手中并不常用，但恰好是本课题的关键。这使课题具备\textbf{不可替代性}——它不是“用 AI 做数学”的又一个泛泛尝试，而是把统计方法论嵌入数学发现流程的一个具体切口。

\paragraph{第二，它满足课程对“体现智能体运行与迭代”的要求。}
第 \ref{sec:roadmap} 节将说明，CRP 循环的每一轮都会产生可记录、可比较、可审计的中间结果：提出的候选、执行的反例搜索、搜索的失败与成功、筛选的决策。这些记录构成完整的运行轨迹（trace），使得“智能体确实在迭代”这件事有据可查，而不必依赖最终是否解决了一个著名猜想。

\paragraph{第三，它的失败也是有信息量的。}
若智能体最终只“重新发现”了若干已知定理，这本身就是对系统正确性的验证；若它提出了在小图上成立、在更大图上被推翻的伪猜想，这些反例恰好量化了“小样本外推”的风险。相比之下，一个以“改进某个数值记录”为目标的课题，一旦数值上不去就几乎没有可汇报的内容。

\begin{notebox}[与课程要求的关系]
课程作业允许不彻底解决原始问题，但要求体现智能体的运行与迭代结果。本课题的“原始问题”（问题 \ref{prob:main}）是一个方法性问题，其价值不依赖于是否发现新的数学定理：即使最终只得到一个“能可靠地重现已知结论、并诚实报告其失败模式”的系统，也已经回答了一个有意义的问题——\emph{这类系统在什么条件下可信}。
\end{notebox}
